\documentclass[tikz,border=8pt]{standalone}
\usepackage{tikz}
\usepackage{amsmath}
\usepackage{amssymb}
\usepackage{pgfplots}
\pgfplotsset{compat=1.18}
\usepackage{ctex}
\usetikzlibrary{calc, positioning, arrows.meta, decorations.pathreplacing}

\begin{document}
\begin{tikzpicture}[>=Stealth]

%% ===== 左图：单位网格（变换前） =====
\begin{scope}[xshift=0cm, yshift=0cm]
  % 单位网格
  \foreach \i in {0,1,2,3} {
    \draw[gray!30, thin] (\i,0) -- (\i,3);
  }
  \foreach \j in {0,1,2,3} {
    \draw[gray!30, thin] (0,\j) -- (3,\j);
  }
  % 坐标轴
  \draw[->, gray!70, thin] (-0.2,0) -- (3.4,0) node[right, font=\scriptsize] {$x$};
  \draw[->, gray!70, thin] (0,-0.2) -- (0,3.4) node[above, font=\scriptsize] {$y$};
  % 单位基向量
  \draw[->, thick, blue!70!black] (0,0) -- (1,0)
        node[below, font=\scriptsize] {$\mathbf{e}_1$};
  \draw[->, thick, red!70!black]  (0,0) -- (0,1)
        node[left,  font=\scriptsize] {$\mathbf{e}_2$};
  % 单位正方形填色
  \fill[blue!8] (0,0) rectangle (1,1);
  \draw[blue!40, thin] (0,0) rectangle (1,1);
  % 标题
  \node[font=\small\bfseries] at (1.5, -0.65) {变换前（单位网格）};
  \node[font=\scriptsize, gray] at (1.5, -1.05) {$I$ 对应的标准坐标系};
\end{scope}

%% ===== 右图：变形网格（A 变换后） =====
\begin{scope}[xshift=7.5cm, yshift=0cm]
  % A = [[1.5, 0.5],[0.5, 1.2]] 的列向量
  % a1 = (1.5, 0.5), a2 = (0.5, 1.2)
  % 变形后的网格线（列向量的整数倍）
  \foreach \i in {0,1,2} {
    % 沿 a1 方向的网格线：从 \i*a2 出发，加 3*a1
    \draw[gray!30, thin]
      ({0.5*\i},{1.2*\i}) -- ({0.5*\i + 3*1.5},{1.2*\i + 3*0.5});
  }
  \foreach \i in {0,1,2} {
    % 沿 a2 方向的网格线：从 \i*a1 出发，加 3*a2
    \draw[gray!30, thin]
      ({1.5*\i},{0.5*\i}) -- ({1.5*\i + 2*0.5},{0.5*\i + 2*1.2});
  }
  % 坐标轴
  \draw[->, gray!70, thin] (-0.2,0) -- (5.2,0) node[right, font=\scriptsize] {$x$};
  \draw[->, gray!70, thin] (0,-0.2) -- (0,3.4) node[above, font=\scriptsize] {$y$};
  % 变形基向量 A*e1 = (1.5,0.5), A*e2 = (0.5,1.2)
  \draw[->, thick, blue!70!black] (0,0) -- (1.5,0.5)
        node[below right, font=\scriptsize] {$A\mathbf{e}_1$};
  \draw[->, thick, red!70!black]  (0,0) -- (0.5,1.2)
        node[left, font=\scriptsize] {$A\mathbf{e}_2$};
  % 变形平行四边形
  \fill[red!8] (0,0) -- (1.5,0.5) -- (2.0,1.7) -- (0.5,1.2) -- cycle;
  \draw[red!50, thin] (0,0) -- (1.5,0.5) -- (2.0,1.7) -- (0.5,1.2) -- cycle;
  % 标题
  \node[font=\small\bfseries] at (2.5, -0.65) {$A$ 变换后（变形网格）};
  \node[font=\scriptsize, gray] at (2.5, -1.05) {$A$ 对应的坐标系};
\end{scope}

%% ===== 箭头 A 和 A^{-1} =====
% 正向箭头：A 变换
\draw[->, very thick, blue!60!black, >=Stealth]
      (3.5, 1.5) -- (7.0, 1.5)
      node[midway, above, font=\small] {$A$（变换）};
% 反向箭头：A^{-1} 还原
\draw[->, very thick, red!60!black, >=Stealth]
      (7.0, 1.0) -- (3.5, 1.0)
      node[midway, below, font=\small] {$A^{-1}$（逆变换）};

%% ===== 底部公式框 =====
\node[font=\small, fill=white, draw=gray!50, thin, rounded corners=2pt,
      inner sep=7pt, align=center] at (6.0, -1.9) {
  $A^{-1}A = I$ \quad（逆变换将网格还原为单位网格）\quad
  $\det(A)\neq 0$ 是可逆的充要条件
};

%% 整体标题
\node[font=\large\bfseries] at (6.0, 4.2) {矩阵逆的几何意义：$A$ 变换与 $A^{-1}$ 反变换};

\end{tikzpicture}
\end{document}
